% コマ10: 2パス構成 — 同じ AST を collect_strings と codegen が順にたどる
% printf("hi"); を含む AST に、目的の違う2本の走査を通す
\documentclass[border=8pt]{standalone}
% 講義資料の図(materials/sessions/figures/*.tex)共通プリアンブル
% 各図の .tex から % 講義資料の図(materials/sessions/figures/*.tex)共通プリアンブル
% 各図の .tex から % 講義資料の図(materials/sessions/figures/*.tex)共通プリアンブル
% 各図の .tex から \input{figure-preamble} で読み込む。
% 配色・フォントは latex/session-template.tex と揃える。

% 日本語(LuaLaTeX + luatexja)。等幅セル内の日本語も和文フォントで出す
\usepackage{luatexja}
\usepackage{luatexja-fontspec}
\setmainfont{Noto Sans CJK JP}
\setmainjfont{Noto Sans CJK JP}
\setmonofont{Inconsolata}[Scale=MatchLowercase]
\setmonojfont{Noto Sans CJK JP}

\usepackage{xcolor}
\definecolor{TcAccentDark}{HTML}{583B66}
\definecolor{TcAccentLight}{HTML}{EEE6F2}
\definecolor{TcAccentLighter}{HTML}{F7F3F9}
\definecolor{TcEmph}{HTML}{E64B17}
\definecolor{TcText}{HTML}{262626}
\definecolor{TcGrayText}{HTML}{737373}
\definecolor{TcGrayBg}{HTML}{F2F2F2}

\usepackage{tikz}
\usetikzlibrary{arrows.meta,positioning,calc,decorations.pathreplacing,fit,backgrounds}

\tikzset{
  % テキスト
  every node/.style={text=TcText},
  annot/.style={font=\small, text=TcText, align=left},
  gannot/.style={font=\small, text=TcGrayText, align=center},
  addr/.style={font=\ttfamily\small, text=TcGrayText},
  % メモリセル
  memcell/.style={draw=TcAccentDark, line width=0.6pt, fill=white,
                  minimum height=8.5mm, minimum width=40mm,
                  font=\ttfamily, inner sep=2pt, outer sep=0pt},
  memcell gray/.style={memcell, fill=TcGrayBg},
  memcell hi/.style={memcell, fill=TcAccentLighter},
  memcell dashed/.style={memcell, dashed, fill=none, text=TcGrayText},
  % 領域(セクション図など)
  region/.style={draw=TcAccentDark, line width=0.6pt, fill=white,
                 minimum width=48mm, align=center, inner sep=3pt, outer sep=0pt},
  % 制御フロー図のボックス
  cfgbox/.style={draw=TcAccentDark, line width=0.6pt, fill=white, rounded corners=1.5mm,
                 minimum width=28mm, minimum height=8mm, font=\small, align=center,
                 inner sep=2pt},
  % 矢印
  ptr/.style={-{Stealth[length=2.8mm]}, line width=1.1pt, color=TcEmph},
  flow/.style={-{Stealth[length=2.4mm]}, line width=0.8pt, color=TcAccentDark},
  flow back/.style={flow, dashed},
  regptr/.style={-{Stealth[length=2.8mm]}, line width=1.1pt, color=TcAccentDark},
  % 補助
  bracestyle/.style={decorate, decoration={brace, amplitude=4pt}, line width=0.8pt,
                     color=TcAccentDark},
}
 で読み込む。
% 配色・フォントは latex/session-template.tex と揃える。

% 日本語(LuaLaTeX + luatexja)。等幅セル内の日本語も和文フォントで出す
\usepackage{luatexja}
\usepackage{luatexja-fontspec}
\setmainfont{Noto Sans CJK JP}
\setmainjfont{Noto Sans CJK JP}
\setmonofont{Inconsolata}[Scale=MatchLowercase]
\setmonojfont{Noto Sans CJK JP}

\usepackage{xcolor}
\definecolor{TcAccentDark}{HTML}{583B66}
\definecolor{TcAccentLight}{HTML}{EEE6F2}
\definecolor{TcAccentLighter}{HTML}{F7F3F9}
\definecolor{TcEmph}{HTML}{E64B17}
\definecolor{TcText}{HTML}{262626}
\definecolor{TcGrayText}{HTML}{737373}
\definecolor{TcGrayBg}{HTML}{F2F2F2}

\usepackage{tikz}
\usetikzlibrary{arrows.meta,positioning,calc,decorations.pathreplacing,fit,backgrounds}

\tikzset{
  % テキスト
  every node/.style={text=TcText},
  annot/.style={font=\small, text=TcText, align=left},
  gannot/.style={font=\small, text=TcGrayText, align=center},
  addr/.style={font=\ttfamily\small, text=TcGrayText},
  % メモリセル
  memcell/.style={draw=TcAccentDark, line width=0.6pt, fill=white,
                  minimum height=8.5mm, minimum width=40mm,
                  font=\ttfamily, inner sep=2pt, outer sep=0pt},
  memcell gray/.style={memcell, fill=TcGrayBg},
  memcell hi/.style={memcell, fill=TcAccentLighter},
  memcell dashed/.style={memcell, dashed, fill=none, text=TcGrayText},
  % 領域(セクション図など)
  region/.style={draw=TcAccentDark, line width=0.6pt, fill=white,
                 minimum width=48mm, align=center, inner sep=3pt, outer sep=0pt},
  % 制御フロー図のボックス
  cfgbox/.style={draw=TcAccentDark, line width=0.6pt, fill=white, rounded corners=1.5mm,
                 minimum width=28mm, minimum height=8mm, font=\small, align=center,
                 inner sep=2pt},
  % 矢印
  ptr/.style={-{Stealth[length=2.8mm]}, line width=1.1pt, color=TcEmph},
  flow/.style={-{Stealth[length=2.4mm]}, line width=0.8pt, color=TcAccentDark},
  flow back/.style={flow, dashed},
  regptr/.style={-{Stealth[length=2.8mm]}, line width=1.1pt, color=TcAccentDark},
  % 補助
  bracestyle/.style={decorate, decoration={brace, amplitude=4pt}, line width=0.8pt,
                     color=TcAccentDark},
}
 で読み込む。
% 配色・フォントは latex/session-template.tex と揃える。

% 日本語(LuaLaTeX + luatexja)。等幅セル内の日本語も和文フォントで出す
\usepackage{luatexja}
\usepackage{luatexja-fontspec}
\setmainfont{Noto Sans CJK JP}
\setmainjfont{Noto Sans CJK JP}
\setmonofont{Inconsolata}[Scale=MatchLowercase]
\setmonojfont{Noto Sans CJK JP}

\usepackage{xcolor}
\definecolor{TcAccentDark}{HTML}{583B66}
\definecolor{TcAccentLight}{HTML}{EEE6F2}
\definecolor{TcAccentLighter}{HTML}{F7F3F9}
\definecolor{TcEmph}{HTML}{E64B17}
\definecolor{TcText}{HTML}{262626}
\definecolor{TcGrayText}{HTML}{737373}
\definecolor{TcGrayBg}{HTML}{F2F2F2}

\usepackage{tikz}
\usetikzlibrary{arrows.meta,positioning,calc,decorations.pathreplacing,fit,backgrounds}

\tikzset{
  % テキスト
  every node/.style={text=TcText},
  annot/.style={font=\small, text=TcText, align=left},
  gannot/.style={font=\small, text=TcGrayText, align=center},
  addr/.style={font=\ttfamily\small, text=TcGrayText},
  % メモリセル
  memcell/.style={draw=TcAccentDark, line width=0.6pt, fill=white,
                  minimum height=8.5mm, minimum width=40mm,
                  font=\ttfamily, inner sep=2pt, outer sep=0pt},
  memcell gray/.style={memcell, fill=TcGrayBg},
  memcell hi/.style={memcell, fill=TcAccentLighter},
  memcell dashed/.style={memcell, dashed, fill=none, text=TcGrayText},
  % 領域(セクション図など)
  region/.style={draw=TcAccentDark, line width=0.6pt, fill=white,
                 minimum width=48mm, align=center, inner sep=3pt, outer sep=0pt},
  % 制御フロー図のボックス
  cfgbox/.style={draw=TcAccentDark, line width=0.6pt, fill=white, rounded corners=1.5mm,
                 minimum width=28mm, minimum height=8mm, font=\small, align=center,
                 inner sep=2pt},
  % 矢印
  ptr/.style={-{Stealth[length=2.8mm]}, line width=1.1pt, color=TcEmph},
  flow/.style={-{Stealth[length=2.4mm]}, line width=0.8pt, color=TcAccentDark},
  flow back/.style={flow, dashed},
  regptr/.style={-{Stealth[length=2.8mm]}, line width=1.1pt, color=TcAccentDark},
  % 補助
  bracestyle/.style={decorate, decoration={brace, amplitude=4pt}, line width=0.8pt,
                     color=TcAccentDark},
}

\begin{document}
\begin{tikzpicture}[node distance=0mm,
    ast/.style={cfgbox, minimum width=30mm, minimum height=7mm, font=\ttfamily\small},
    p1/.style={flow, line width=1.1pt},
    p2/.style={ptr},
    num/.style={annot, font=\small\bfseries, text=TcEmph, align=center},
    numa/.style={annot, font=\small\bfseries, text=TcAccentDark, align=center},
    note/.style={annot, text=TcGrayText},
    lcline/.style={font=\ttfamily\footnotesize, text=TcText, anchor=west}]

  \node[gannot, anchor=north] at (0mm, 46mm)
      {{\ttfamily main() \{ printf("hi"); \}} の AST（1 本だけ）を、目的の違う 2 つの走査が順に通る};

  % ================ 中央: AST ================
  \node[ast] (stmt) at (0mm, 20mm)  {ExprStmt};
  \node[ast] (call) at (0mm, 4mm)   {Call printf};
  \node[ast] (str)  at (0mm, -12mm) {Str "hi"};
  \draw[flow] (stmt.south) -- (call.north);
  \draw[flow] (call.south) -- (str.north);

  % ================ パス1（左・collect_strings） ================
  \node[numa, anchor=south] at (-38mm, 29mm)
      {① パス1\\{\ttfamily collect\_strings\_stmt()}};
  \draw[p1] (-38mm, 26mm) -- (-38mm, -8mm);
  \node[note, anchor=east, text=TcAccentDark, align=right] at (-46mm, 8mm)
      {木を上から順にたどり、\\{\ttfamily Str} に着いたら\\{\ttfamily \_intern()} で登録する};

  \node[region, minimum width=64mm, minimum height=13mm,
        font=\ttfamily\footnotesize, align=center] (dict) at (-64mm, -34mm)
      {self.\_strings = \{"hi": ".LC1"\}};
  \draw[p1] (str.south) to[out=-135, in=90] (dict.north);

  % ================ パス2（右・codegen） ================
  \node[num, anchor=south] at (38mm, 29mm)
      {② パス2\\{\ttfamily gen\_func()}};
  \draw[p2] (38mm, 26mm) -- (38mm, -8mm);
  \node[note, anchor=west, text=TcEmph, align=left] at (46mm, 8mm)
      {同じ木をもう一度たどり、\\{\ttfamily Str} では登録済みの\\ラベルを命令に埋める};

  \node[region, minimum width=44mm, minimum height=13mm,
        font=\ttfamily\footnotesize, align=center] (insn) at (64mm, -34mm)
      {la a0, .LC1};
  \draw[p2] (str.south) to[out=-45, in=90] (insn.north);

  % ================ 下: 出力されるアセンブリ ================
  \node[gannot, anchor=south] at (0mm, -49mm) {出力されるアセンブリ（{\ttfamily .data} が先、{\ttfamily .text} が後）};

  \node[region, minimum width=108mm, minimum height=17mm, anchor=north] (data) at (0mm, -51mm) {};
  \node[lcline] at ($(data.north west)+(4mm,-4mm)$)  {.data};
  \node[lcline] at ($(data.north west)+(4mm,-9mm)$)  {.LC1:};
  \node[lcline] at ($(data.north west)+(4mm,-14mm)$) {\ \ .byte 104, 105, 0};
  \node[numa, anchor=east] at ($(data.east)+(-4mm,0mm)$) {パス1 が埋める};

  \node[region, minimum width=108mm, minimum height=17mm, anchor=north] (text) at (0mm, -71mm) {};
  \node[lcline] at ($(text.north west)+(4mm,-4mm)$)  {.text};
  \node[lcline] at ($(text.north west)+(4mm,-9mm)$)  {main:};
  \node[lcline] at ($(text.north west)+(4mm,-14mm)$) {\ \ la a0, .LC1};
  \node[num, anchor=east] at ($(text.north east)+(-4mm,-5mm)$) {パス2 が埋める};

  \draw[p1] (dict.south) to[out=-90, in=180] ($(data.west)+(0mm,-3mm)$);
  \draw[p2] (insn.south) to[out=-90, in=0]  ($(text.east)+(0mm,-4mm)$);

  % ---- 補足 ----
  \node[gannot, anchor=north, align=center] at (0mm, -94mm)
      {{\ttfamily .data} を先に出すには、コードを書き始める前に文字列を全部知っている必要がある。\\
       だからパス1 が木を最後までたどり終えるまで、パス2 は始められない};
\end{tikzpicture}
\end{document}
